% !PRINTED 83
\weberpage{83}

Wir machen nun die Annahme, dass unter den $\nu$-reihigen Determinanten der Matrix \gesperrt{wenigstens eine} von Null verschieden sei, während die $(\nu+1)$reihigen und folglich auch die höheren Determinanten, falls solche vorhanden sind, alle verschwinden sollen. $\nu$ kann jede Zahl sein, die nicht grösser als die kleinere der beiden Zahlen $n$ oder $m$ ist (oder falls $n=m$ ist, diesen gemeinschaftlichen Werth nicht übertrifft).

Eine solche Zahl $\nu$ wird sich immer finden lassen, wenn wir den schon erledigten, ganz interesselosen Fall ausschliessen, dass alle Coëfficienten $a_i^{(k)}$ verschwinden.

Wir können, ohne die Allgemeinheit zu beschränken, zur Vereinfachung der Bezeichnung annehmen, die nicht verschwindende $\nu$-reihige Determinante sei
\begin{equation}
 A=\begin{vmatrix}
 a_1^{(1)}&a_2^{(1)}&\cdots&a_\nu^{(1)}\\
 a_1^{(2)}&a_2^{(2)}&\cdots&a_\nu^{(2)}\\
 \vdots&\vdots&&\vdots\\
 a_1^{(\nu)}&a_2^{(\nu)}&\cdots&a_\nu^{(\nu)}
 \end{vmatrix}.
\tag{4}
\end{equation}
Denn offenbar steht es uns frei, das Gleichungssystem (1) in beliebiger Weise anzuordnen, und ferner können wir die Bezeichnung der Unbekannten $x$ so wählen, dass irgend $\nu$ von ihnen die $\nu$ ersten sind. Die Unterdeterminanten von $A$ bezeichnen wir wie früher mit $A_i^{(k)}$, worin $i,k$ von $1$ bis $\nu$ gehen.

\gesperrt{I. Wenn nun zunächst $\nu=n$ ist, was voraussetzt, dass $m$ nicht kleiner als $n$ ist, so haben die Gleichungen (1) keine andere Lösung, als die in den Gleichungen (2) enthaltene.}

Denn greifen wir die $n$ ersten der Gleichungen (1) heraus:
\begin{equation}
\begin{aligned}
 a_1^{(1)}x_1+a_2^{(1)}x_2+\cdots+a_n^{(1)}x_n&=0,\\
 a_1^{(2)}x_1+a_2^{(2)}x_2+\cdots+a_n^{(2)}x_n&=0,\\
 &\vdots\\
 a_1^{(n)}x_1+a_2^{(n)}x_2+\cdots+a_n^{(n)}x_n&=0,
\end{aligned}
\tag{5}
\end{equation}
und multipliciren diese der Reihe nach mit $A_\mu^{(1)},A_\mu^{(2)},\ldots,A_\mu^{(n)}$, worin $\mu$ jeder der Indices $1,2,\ldots,n$ sein kann, und addiren sie, so folgt, weil nach \S.\,22 (2) und (6)
\[
 \sum_{1,\nu}^{i} A_\mu^{(i)} a_\lambda^{(i)}=0 \quad\hbox{oder}\quad =A
\]
%CONT>>
